\section{Related Work}
\label{sec:related-work}

\subsection{Circuit simulators}

SPICE and its descendants~\cite{nagel1973spice} remain the reference implementation of
circuit simulation and are built around modified nodal analysis~\cite{ho1975mna}: node
voltages and a small number of auxiliary branch currents (one per independent voltage source
or inductor, in the classical formulation) are solved simultaneously at each timestep, with
reactive elements replaced by companion resistive/source models derived from a numerical
integration rule. Educational tools such as the Falstad circuit simulator and PhET's Circuit
Construction Kit lower the barrier to entry by replacing SPICE's netlist syntax with direct
schematic manipulation and are widely used in introductory courses. All three tools,
however, present the circuit as a two-dimensional schematic. CircuitCraft uses the same
underlying numerical technique (Sections~\ref{sec:architecture}
and~\ref{sec:circuit-solver} describe an MNA solver employing trapezoidal-integration
companion models for reactive elements, the family of technique SPICE itself uses) but
embeds it in a three-dimensional, first-person, block-based world rather than a schematic
canvas.

\subsection{Minecraft in education}
\label{sec:minecraft-in-education}

The instructional value of games and simulations has been studied extensively in general: a
meta-analysis by Sitzmann~\cite{sitzmann2011meta} found simulation games to be more
effective for training self-efficacy and knowledge retention than comparison instructional
methods across a range of technical domains. Minecraft occupies a disproportionately large
share of that broader literature. Microsoft maintains and markets a dedicated,
classroom-oriented variant of the game, \emph{Minecraft: Education Edition}, an indication of
how extensively the game had already been adopted informally by educators prior to this
commercial productization. Nebel et al.'s review~\cite{nebel2016mining} surveys this
literature broadly and identifies Minecraft's principal appeal for education as the same
property that makes it appealing as a game: an open, constructive world in which a learner's
own artifact, rather than a worked example, is the object of study, closely matching
Papert's constructionist framing~\cite{papert1980mindstorms}, discussed further in
Section~\ref{sec:pedagogy}.

Published case studies span a range of subjects well outside engineering. Bos et
al.~\cite{bos2014mathematics} describe using Minecraft to teach elementary-school
mathematics concepts, including area, perimeter, and ratio, through building tasks; Short's
account~\cite{short2012teaching} covers the use of the game to teach general science concepts
in a secondary-school setting; and Cipollone et al.~\cite{cipollone2014creative} present a
case study of Minecraft used as an open-ended creative tool, emphasizing collaboration and
planning skills rather than any specific curricular content. Callaghan's
survey~\cite{callaghan2016investigating} of Minecraft's role across primary and secondary
classrooms identifies a theme common to each of the studies above: teachers report high
engagement and motivation fairly consistently, whereas rigorously measured gains in specific
learning outcomes are reported less often and less consistently, a caveat also present in
the general simulation-games literature~\cite{sitzmann2011meta} and one this paper does not
claim to have resolved for CircuitCraft specifically (Section~\ref{sec:limitations}).

Within that broader literature, the subset closest to the present paper's subject concerns
Minecraft's native redstone subsystem, already a common entry point for teaching digital
logic and computer architecture concepts: players build logic gates, adders, and working
memory cells entirely from in-game blocks. Dezuanni et al.~\cite{dezuanni2015redstone}
document children describing and reasoning about redstone circuits in explicitly electrical
language, notwithstanding that the underlying mechanic is a discrete, boolean signal rather
than a continuous one (the mechanics of redstone itself are described in
Section~\ref{sec:minecraft-mechanics}). This is the literature to which CircuitCraft is
closest, and from which it departs most directly: the case studies cited above use
Minecraft's existing mechanics, redstone included, as the vehicle for a lesson, whereas Mine
Memristors introduces an entirely new, continuous-time simulation layer, because redstone's
discrete signal cannot represent voltage, current, or reactive state at all.

Beyond that academic literature, several existing Minecraft \emph{mods} -- as opposed to
case studies built on the base game -- add circuit-building mechanics of their own, and are
worth situating CircuitCraft against directly. Create: Circuits~\cite{chocoboyy2026createcircuits}
and ProjectRed~\cite{mrtjp2026projectred} both extend the redstone subsystem with more
compact logic-gate, arithmetic, and memory-latch blocks than vanilla redstone alone provides
-- ProjectRed as a modular suite (its Integration module specifically supplies the logic
gates, alongside separate modules for wiring, fabrication, and other redstone-adjacent
concerns). Create: Circuits labels some of its blocks (adder, multiplier, min, max) as
handling ``analog'' values, but only in the sense of Create's own multi-valued signal
strength, an integer quantity in the same fundamentally discrete sense as ordinary redstone
(Section~\ref{sec:minecraft-mechanics}), not a continuous electrical one; like ProjectRed, it
represents no voltage, current, or reactive state whatsoever, and both mods are purely
logical extensions of the same subsystem discussed above rather than a departure from it.
CircuitSim~\cite{dankov2026circuitsim} is, to the author's knowledge, the only prior
Minecraft mod that simulates genuinely analog circuits -- resistors, capacitors, inductors,
diodes, transistors, and op-amps among its components. It does so, however, by exporting a
player's in-world circuit as a netlist and invoking an external ngspice~\cite{nagel1973spice}
process, installed separately on the host system and on the system \texttt{PATH}, each time a
player right-clicks a dedicated ``simulate'' block, rather than solving the circuit natively
inside the mod itself. CircuitCraft's solver, by contrast, has no external dependency of any
kind, is embedded directly in the mod, and runs every game tick as an intrinsic part of the
running world -- driving the oscilloscope HUDs and the memristor's own evolving state live,
rather than producing a report a player consults after invoking a separate tool. To the
author's knowledge, it is the first Minecraft modification to embed a native, continuous-time
circuit solver for this purpose, rather than approximating one with the game's existing
discrete redstone mechanics or delegating it to an external simulator.

\subsection{Memristor theory and compact models}

The memristor was postulated on symmetry grounds by Chua~\cite{chua1971memristor} as a
fourth fundamental two-terminal circuit element, relating charge and flux, alongside the
resistor, capacitor, and inductor; a physical device exhibiting the predicted
charge-controlled behavior was subsequently reported by Strukov et al.\ at
HP~Labs~\cite{strukov2008missing}, whose linear-drift formulation is the model implemented
in this mod (Section~\ref{sec:circuit-solver}).

A substantial body of subsequent work refines or replaces this original formulation,
targeting shortcomings that only become apparent once the model is used for large-scale
circuit design or compared against fabricated devices. Biolek et
al.~\cite{biolek2009spice} introduce a window function that prevents the state variable
from becoming permanently pinned at either boundary during SPICE simulation, a numerical
issue that the original linear-drift formulation is otherwise prone to; Kvatinsky et
al.'s TEAM model~\cite{kvatinsky2013team} instead adopts a threshold-based switching
condition, trading some physical detail for a form that is considerably cheaper to simulate
at the scale of a large memristor-based circuit. Guan, Yu, and Wong~\cite{guan2012compact}
take a converse approach, developing a physics-based compact model of filament formation and
rupture, validated directly against measured resistive-switching device data, at the cost of
substantially greater model complexity. Picos et al.~\cite{picos2015chargeflux} formulate an
alternative approach that works directly in the charge--flux domain rather than through an
explicit resistance-state variable, and Al Chawa, Picos, and
Tetzlaff~\cite{alchawa2020simple} present a simplified compact model aimed specifically at
neuromorphic ReRAM applications, where computational efficiency across large device arrays
is often prioritized over matching any single device's measured characteristics exactly.

CircuitCraft implements Strukov et al.'s original linear-drift formulation specifically,
rather than any of the refinements above, because the priorities of a teaching tool differ
from those of a circuit designer or device physicist: the model must be simple enough to
verify in closed form (Section~\ref{sec:validation}) and traceable, step by step, to the
historically foundational definition a student is most likely to encounter first, rather
than optimized for matching a specific fabricated device's current--voltage characteristic
or for simulation efficiency across a large memristor array. Exposing one or more of the
alternative models above as a selectable preset -- in particular a windowed variant such as
Biolek et al.'s, which addresses a real limitation of the current implementation discussed
in Section~\ref{sec:limitations} -- is a natural extension of the present work.
